\chapter{選擇計算引數}\label{ch:theory}

完成一次計算之後，接下來的問題通常是：這些引數為什麼這樣選，換一個體系時應該改哪些，以及怎樣判斷精度已經足夠。本章圍繞介面中的實際設定介紹電子結構計算的基本概念。閱讀時可以保留第\ref{ch:first}章的氫分子專案，一邊檢視引數，一邊理解它們所描述的物理問題。

一套完整的計算條件包括幾何結構、電子模型和數值精度。幾何結構給出原子種類、位置與晶格；電子模型給出泛函、贗勢和自旋等假設；數值精度則涉及截斷能、k 點取樣和迭代收斂條件。三者各有作用。數值上充分收斂的計算，仍然是在所選模型下求得的答案；模型選得合適，也需要足夠的數值精度才能得到穩定結果。

\section{從原子位置到電子密度}

\subsection{密度泛函理論解決什麼問題}
\index{Midu fanhan lilun@密度泛函理論}

在給定原子位置時，電子並不是靜止在原子周圍的一組小球。它們的運動和相互作用共同決定體系的能量、化學鍵和許多材料性質。直接追蹤多電子波函式十分困難，密度泛函理論改用空間中的電子密度來描述基態問題。

電子密度可以理解為電子在空間中分佈的程度。密度較大的區域通常有較多電子，但它並不等於檢視器中原子球的大小。球棍模型幫助我們辨認幾何結構，電子密度則來自計算；兩者的影象形態和用途不同。

KSSOLV 透過 Kohn--Sham 方程求解電子態。用於求解的有效勢取決於電子密度，而電子密度又由電子態產生，因此不能只計算一次便結束。求解器需要不斷更新密度，直到方程所用的密度與求得的密度互相一致，這就是自洽。


在本書非磁性、PBE贗勢教程中，可以把所求的方程寫為
\begin{equation}
  \begin{aligned}
  \widehat H[\rho]\psi_{n\mathbf k}
    &=\varepsilon_{n\mathbf k}\psi_{n\mathbf k},\\
  \widehat H[\rho]&= -\frac12\nabla^2
  +V_{\mathrm{loc}}^{\mathrm{ps}}+\widehat V_{\mathrm{NL}}^{\mathrm{ps}}
  +V_{\mathrm H}[\rho]+V_{\mathrm{xc}}[\rho].
  \end{aligned}
  \label{eq:ks-operator}
\end{equation}
這裡採用原子單位：長度為Bohr，能量為Ha。動能項描述電子的空間變化；局域和非局域贗勢描述原子實；Hartree勢與交換關聯勢隨密度更新。介面中的\ui{Pseudopotential}和\ui{Functional}分別決定相應部分。HSE06還含非局域交換處理，不能僅用上式中的局域交換關聯勢概括。

非磁性設定\ui{Nspin}=1時，電子密度由這些態加權得到：
\begin{equation}
  \begin{aligned}
  \rho(\mathbf r)&=2\sum_{\mathbf k}w_{\mathbf k}
       \sum_n f_{n\mathbf k}|\psi_{n\mathbf k}(\mathbf r)|^2,\\
  N_{\mathrm e}&=\int_{\Omega}\rho(\mathbf r)\,\mathrm d^3r.
  \end{aligned}
  \label{eq:ks-density}
\end{equation}
波函式在一個計算晶胞內歸一化。前面的2計入非磁性自旋簡併；滿佔據的一個空間態取$f=1$，對應兩個電子，不應再把$f$寫成2。分子教程只有一個$\Gamma$點，其權重為1。贗勢計算中的$N_{\mathrm e}$是實際處理的價電子數，密度使用電子數每單位體積表示；它與含負電子電荷的電荷密度有符號上的區別。

\begin{table}[htbp]
\centering
\caption{閱讀電子結構公式時的常用符號。}
\begin{tabularx}{\linewidth}{lX}
\toprule
符號 & 含義及與設定的聯絡\\
\midrule
$\psi_{n\mathbf k}$、$\varepsilon_{n\mathbf k}$ & 第$n$個電子態及其能量；$\mathbf k$標記晶體中的波矢\\
$\rho$、$\Omega$ & 價電子密度及計算晶胞的體積；原子單位下分別用$\mathrm{Bohr}^{-3}$和$\mathrm{Bohr}^{3}$\\
$f_{n\mathbf k}$ & 佔據權重，由電子數、\ui{Smearing}和\ui{Temperature}共同確定\\
$w_{\mathbf k}$ & $k$點積分權重；一個自旋通道內的權重和為1\\
\bottomrule
\end{tabularx}
\end{table}

改變\ui{Extra bands}可以增加參與求解的空態，但不會改變式\eqref{eq:ks-density}所要求的電子數。選擇自旋分辨或非共線模式時，需要分別處理相應通道，不能直接沿用這裡的簡併因子2。

\subsection{為什麼固定原子也要迭代}

可以把一次 SCF 看成反覆改進一份電子分佈的過程。首先建立一個初始密度，用它構造有效勢；隨後求解電子態，根據這些態和電子佔據數得到新密度；再把新舊資訊適當混合，作為下一步的起點。

\begin{figure}[htbp]
\centering
\begin{tikzpicture}[x=1mm,y=1mm,font=\sffamily\small,
  stage/.style={draw=teal,fill=paleblue,rounded corners=1.5pt,
    minimum width=31mm,minimum height=11mm,align=center,text=brand},
  flow/.style={->,>=stealth,line width=.7pt,draw=teal}]
  \node[stage] (input) at (15,32) {输入密度\\$\rho_{\mathrm{in}}$};
  \node[stage] (potential) at (55,32) {构造有效势};
  \node[stage] (states) at (95,32) {求电子态与占据};
  \node[stage] (output) at (95,11) {得到新密度\\$\rho_{\mathrm{out}}$};
  \node[stage] (check) at (55,11) {达到收敛条件？};
  \node[stage] (mix) at (15,11) {混合新旧密度};
  \node[text=brand,align=center] (done) at (55,-6) {保留收敛结果};
  \draw[flow] (input) -- (potential);
  \draw[flow] (potential) -- (states);
  \draw[flow] (states) -- (output);
  \draw[flow] (output) -- (check);
  \draw[flow] (check) -- node[above,text=brand] {否} (mix);
  \draw[flow] (mix) -- (input);
  \draw[flow] (check) -- node[right,text=brand] {是} (done);
\end{tikzpicture}
\caption{电子自洽循环（示意）。在原子位置不变的条件下更新电子密度，直到输入与输出满足收敛要求。}
\label{fig:scf-cycle}
\end{figure}


這個過程中，原子位置可以完全不變。因此，SCF 的“第十步”不表示原子移動了十次。結構最佳化還包含另一層迴圈：先在一個幾何構型上完成電子自洽，再根據原子力移動原子，然後對新構型重新計算電子結構。區分電子迭代和原子位置迭代，是閱讀日誌的重要基礎。

\subsection{自洽與非自洽計算的分工}
\index{Ziqiachang@自洽場（SCF）}

SCF 的目標是建立與結構和模型一致的電子密度。能帶和態密度還需要在更多或不同的 k 點上求電子能量，這時可以使用已經收斂的密度，保持有效勢不再隨這些附加取樣反覆自洽。這類計算通常稱為非自洽計算，簡稱 NSCF。

在工作流中，\ui{SCF} 先提供收斂密度；\ui{NSCF} 指定後續求解所需的演算法選項；\ui{Band Calculation} 或 \ui{DOS} 再根據相應的路徑或網格執行計算。能帶路徑和 DOS 網格描述的是不同的取樣任務，不應因為它們都含有 k 點，就把兩者設成同一個列表。

\section{選擇分子還是晶體}

\subsection{計算物件的含義}

\ui{Build Molecule} 中的 \ui{Molecule} 與 \ui{Crystal} 決定如何組織計算物件。對於孤立分子教程，選擇 Molecule，併為分子保留合適的模擬盒；對於週期性材料，選擇 Crystal，並給出晶格和 k 點取樣。

“分子”並不意味著計算完全沒有盒子。平面波方法仍然需要一個空間區域來展開波函式。盒內分子與其週期重複像之間的相互作用，應透過足夠的真空距離減小。對於帶電、強極性或空間分佈很寬的體系，所需盒尺寸可能與中性緊湊分子不同。

分子檔案、晶體檔案與計算型別也不能完全等同。CIF 通常儲存晶體，但也可以儲存分子示例的模擬盒和其中的分子座標。XYZ 通常儲存分子座標，匯入後軟體可以為它構造一個盒子。最終仍應檢查實際幾何資訊，並在節點中明確選擇計算型別。

\subsection{晶格向量與原子座標}

週期結構由三個晶格向量和晶胞內的原子位置組成。晶格的長度、夾角決定晶胞形狀；原子位置可以用直角座標表示，也可以用沿三個晶格向量的分數座標表示。

分數座標 $(0.25,0.25,0.25)$ 表示沿三個晶格向量各取四分之一後相加，並不總等於三個直角座標都為晶格常數的四分之一。只有在相應的正交晶胞中，二者才有這種簡單關係。對於六方、單斜等晶胞，輸入座標時尤其應確認當前座標型別。


用$\mathbf a_i$表示三個實空間晶格向量，$f_i$表示分數座標，則位置與倒格向量滿足
\begin{equation}
  \begin{gathered}
  \mathbf r=\sum_{i=1}^{3}f_i\mathbf a_i,
  \qquad \mathbf k=\sum_{j=1}^{3}\kappa_j\mathbf b_j,\\
  \mathbf a_i\!\cdot\!\mathbf b_j=2\pi\delta_{ij}.
  \end{gathered}
  \label{eq:reciprocal-convention}
\end{equation}
$\delta_{ij}$在$i=j$時為1，否則為0。這裡的倒格約定包含$2\pi$；$\kappa_j$是倒格基底下的無量綱座標，而$\mathbf k$本身具有長度倒數的單位。因而路徑中的$(0.5,0,0)$表示沿第一倒格向量取一半，並非實空間中的半個晶格長度。

一個晶體可以使用原胞、常規晶胞或超胞表示。它們可能描述同一種無限週期結構，卻包含不同數量的原子，也有不同大小的布里淵區。超胞越大，單位晶胞內的電子態越多；同樣的整數 k 點網格不再代表相同的取樣密度。

\subsection{單位的三個層次}

結構編輯和常用檔案格式通常使用 Å，數值計算中的幾何量使用 Bohr，電子能量則可能在不同位置以 Ha、Ry 或 eV 表示。軟體會在匯入和求解之間完成必要的轉換，使用者仍需在讀取外部資料或整理結果時關注單位。

\begin{table}[htbp]
\centering
\caption{本書常見量與單位}
\begin{tabularx}{\linewidth}{lX}
\toprule
量 & 常見表示與閱讀位置\\
\midrule
幾何長度 & 檢視器、CIF 和通常的 XYZ 用 Å；數值介面可使用 Bohr\\
截斷能 & Build Molecule 的 Ecut 用 Hartree，即 Ha\\
總能量 & 收斂曲線左軸用 Ha；日誌也可能同時列出 Ry\\
能帶與 DOS 能量 & eV，本書教程以費米能為參考零點\\
電子溫度 & Temperature 用 K\\
原子力 & 數值計算記錄採用原子單位，比較外部資料時應換算\\
\bottomrule
\end{tabularx}
\end{table}

常用換算為 $1\,\mathrm{Ha}=2\,\mathrm{Ry}\approx27.2114\,\mathrm{eV}$，$1\,\mathrm{Bohr}\approx0.529177\,\text{Å}$。如果只是檢查輸入，使用這些近似值即可；批次整理數值時，應使用程式中的轉換常數，避免在不同表格中混用過多種舍入值。

\section{泛函：選擇電子相互作用的近似}

\subsection{常用泛函}

\ui{Functional} 提供 PBE、PZ 和 HSE06。它們採用不同方式近似電子的交換關聯能。PZ 屬於局域密度近似一類，PBE 屬於廣義梯度近似，HSE06 屬於遮蔽雜化泛函。它們的區別影響計算得到的平衡結構、相對能量和電子能級。

對於學習工作流和一般的初步材料計算，PBE 是便於開始的選擇。本書的基礎教程統一使用 PBE，使讀者能夠把注意力放在結構和計算過程上。切換到其他泛函時，不能只比較執行時間，還應考慮贗勢是否合適、研究目標是什麼，以及需要比較的文獻採用了哪些條件。

HSE06 引入了一部分精確交換，計算量通常高於半局域泛函。它還需要相應的交換演算法和後續計算條件。初次學習時，先完成 PBE 的 SCF、能帶和 DOS，再用相同的小體系學習雜化泛函的演算法設定。

\subsection{怎樣比較不同泛函}

比較泛函時，應先決定比較的是“同一結構上的電子模型差異”，還是“各自平衡結構上的結果差異”。前者保持幾何不變，只改變電子模型；後者允許不同模型得到不同的最穩定結構。這兩種比較回答的問題不同。

例如，在固定鍵長上比較兩個氫分子總能量，並不能直接判斷哪種泛函給出的平衡鍵長更接近實驗。若研究目標是鍵長，應該分別最佳化結構；若目標是固定結構下的電子能級，則應使用相同幾何。寫計算說明時把這個選擇說明白，結果會更容易理解。

\section{贗勢：怎樣描述原子實}
\index{Yanshi@贗勢}

贗勢把緊密束縛的內層電子和原子核的影響合併為有效描述，使計算主要處理價電子。這是平面波計算能夠高效處理多種元素的重要原因。贗勢本身包含建模選擇，因此它也是計算條件的一部分。

\ui{Pseudopotential} 的 \ui{default} 使用隨軟體提供的預設贗勢集合。本書的氫和矽教程使用內建 PBE-1.3 ONCV 檔案。在執行訊息中可以檢視實際載入的檔名稱，確認每一種元素都找到了對應的贗勢。

泛函和贗勢應配合使用。更改 Functional 不表示軟體會自動選擇一套新的贗勢，也不表示所有組合都同樣合適。開始新的材料研究時，應記錄贗勢來源、版本和價電子選擇，並在同一組可比較的計算中保持一致。

\subsection{為什麼不同贗勢需要不同截斷能}
\index{Jieduanneng@截斷能}

有的贗勢產生較平滑的價電子波函式，較少的平面波便能描述；有的需要更高的截斷能。截斷能是否足夠，取決於體系和所用贗勢，而不是隻取決於元素的原子序數或原子數量。

當一個結構含有多種元素時，應針對完整體系進行檢查。某個元素對平面波精度要求較高，往往會決定整個計算需要的截斷能。氧化物、過渡金屬或含氫體系應分別進行截斷能收斂檢查。

\section{平面波截斷能}
\index{Pingmianbo@平面波}

\subsection{\headinglatin{Ecut} 控制什麼}

平面波把電子波函式表示為空間波的疊加。波長越短，能夠描述越細的空間變化，相應的動能越高。截斷能規定保留哪些平面波，從而控制表示電子態的精細程度。


晶體中的一個電子態可以展開為
\begin{equation}
  \begin{aligned}
  \psi_{n\mathbf k}(\mathbf r)
   &=\frac{1}{\sqrt{\Omega}}\sum_{\mathbf G}
      c_{n\mathbf k}(\mathbf G)
      \exp\!\left[\mathrm i(\mathbf k+\mathbf G)\cdot\mathbf r\right],\\
  T_{\mathbf k+\mathbf G}&=\frac{|\mathbf k+\mathbf G|^2}{2}.
  \end{aligned}
  \label{eq:plane-wave-expansion}
\end{equation}
$\mathbf G$是倒格向量的整數線性組合，$c$是待求係數；採用Bohr及其倒數時，右側動能的單位為Ha。對$\Gamma$點，$\mathbf k=0$，基礎平面波集合滿足$|\mathbf G|^2/2\leq E_{\mathrm{cut}}$。提高截斷能，就是把更遠處、波長更短的分量納入表示。

\begin{figure}[htbp]
\centering
\begin{tikzpicture}[x=1mm,y=1mm,font=\sffamily\small]
  \node[text=brand] at (26,47) {$\Gamma$点的基础平面波集合};
  \node[text=brand] at (88,47) {波长与空间细节};
  \fill[paleblue] (26,22) circle (18);
  \foreach \x in {-24,-18,...,24}
    \foreach \y in {-18,-12,...,18} {
      \pgfmathtruncatemacro{\inside}{\x*\x+\y*\y<=324}
      \ifnum\inside=1
        \fill[brand] ({26+\x},{22+\y}) circle (.75);
      \else
        \fill[muted!35] ({26+\x},{22+\y}) circle (.55);
      \fi
    }
  \draw[teal,line width=.8pt] (26,22) circle (18);
  \draw[->,draw=muted,line width=.4pt] (1,22) -- (53,22);
  \draw[->,draw=muted,line width=.4pt] (26,0) -- (26,43);
  \node[anchor=west,text=muted] at (51,19) {$G_x$};
  \node[anchor=west,text=muted] at (27,42) {$G_y$};
  \node[fill=paleblue,inner sep=1pt,text=brand] at (26,22) {$\Gamma$};
  \draw[draw=muted!40,line width=.4pt] (66,31) -- (112,31);
  \draw[draw=muted!40,line width=.4pt] (66,12) -- (112,12);
  \draw[teal,line width=.8pt,domain=0:44,samples=100]
    plot ({67+\x},{31+3*sin(\x*16.3636)});
  \draw[brand,line width=.8pt,domain=0:44,samples=140]
    plot ({67+\x},{12+3*sin(\x*49.0909)});
  \node[text=ink] at (89,39) {较长波长};
  \node[text=ink] at (89,20) {较短波长};
  \node[text=ink] at (26,-7) {提高Ecut，保留更多分量};
  \node[text=ink] at (89,-7) {短波分量需要更高动能};
\end{tikzpicture}
\caption{平面波截断的二维示意。左图蓝点为截断圆内保留的倒格分量；实际计算使用三维集合。右图说明波长与空间变化的关系。}
\label{fig:plane-wave-cutoff-concept}
\end{figure}


非零$k$點的動能取決於$\mathbf k+\mathbf G$。晶體計算可以共用按基礎$\mathbf G$選擇的集合，因此圖中的圓應理解為$\Gamma$點示意，而不是每個$k$點都重新裁切出完全相同的動能球。實際使用時仍透過收斂檢查判斷\ui{Ecut}是否足夠。

提高 \ui{Ecut (Hartree)} 通常會增加平面波數、記憶體需求和每一步計算時間。它改善的是基組描述，並不會自動增加 k 點取樣，也不會自動放寬或收緊電子自洽條件。幾個精度引數之間相互影響，但不能互相替代。

\subsection{怎樣做截斷能收斂檢查}

先固定結構、盒子、泛函、贗勢和 k 點，再依次提高 Ecut。例如，為一個小型示例體系計算 $15$、$20$、$25$、$30$ Ha 四個值，記錄每次收斂後的能量和所關心的性質。這裡的序列只是檢查方法的示例，最終範圍應根據實際體系調整。

判斷時關注相鄰設定之間的變化，而不是某個總能量看起來是否“足夠負”。如果研究的是相對能量，可以記錄每原子能量差或兩個結構之間的能量差；如果研究的是幾何最佳化，應同時檢查力和最佳化後的尺寸。

\begin{steps}
\item 從同一模板另建一條工作流，或先儲存已有工作流為 \file{.wf}，再透過從檔案匯入建立副本，保留相同結構。
\item 每次只修改 Ecut，並讓 SCF 達到同樣的收斂條件。
\item 把結果整理為截斷能與目標性質的表格。
\item 選擇在所需精度範圍內已經穩定的設定，再進行正式計算。
\end{steps}

SCF 未收斂的結果不宜混入這樣的比較。否則觀察到的差異可能來自迭代尚未結束，而不是真正的基組精度變化。收斂檢查通常比盲目選擇很大的引數更節省整體工作量。

\section{k 點與布里淵區取樣}
\index{K dian caiyang@k 點取樣}

\subsection{為什麼晶體需要 k 點}

週期結構中的電子態可以用波矢 k 來標記。許多總能量和密度有關的量需要對整個布里淵區求和或積分，實際計算則用有限數量的 k 點近似。\mbox{\ui{Autokpts}} 指定規則網格，前三個整數給出三個方向上的網格數，後三項指定相應的網格位移。


對布里淵區內某個量$Q(\mathbf k)$，區域平均可用加權和近似：
\begin{equation}
  \langle Q\rangle_{\mathrm{BZ}}
      \approx\sum_{\mathbf k}w_{\mathbf k}Q(\mathbf k),
  \qquad \sum_{\mathbf k}w_{\mathbf k}=1.
  \label{eq:k-point-average}
\end{equation}
規則完整網格可使用等權重。增加網格密度，改善的是這類區域取樣；沿一條路徑加入更多點，只能使路徑上的曲線更細，不能代替整個布里淵區的積分網格。

例如 \mbox{\texttt{[2 2 2 0 0 0]}} 描述本書矽教程中的起始網格。它是方便學習的小網格，不是矽的所有性質都已充分收斂的結論。增加到更密網格後，可以檢查總能量、力或能帶邊緣是否已經穩定。

對於細長或扁平的晶胞，不同方向不一定需要相同的整數。實空間週期越長，相應倒空間方向的區間越短，所需取樣數通常也會改變。應結合晶胞形狀和目標精度選擇網格，而不是隻追求三個數字相同。

\subsection{為任務選擇取樣方式}

SCF 網格用於建立電子密度和有效勢。能帶路徑選擇若干具有對稱意義的方向，展示能量如何隨 k 變化。DOS 網格則覆蓋整個布里淵區，用於統計一定能量附近有多少電子態。

\begin{figure}[htbp]
\centering
\begin{tikzpicture}[x=1mm,y=1mm,font=\sffamily\small]
  \node[text=brand] at (21,43) {SCF 与 DOS 网格};
  \node[text=brand] at (85,43) {能带路径};
  \fill[paleblue] (3,0) rectangle (39,36);
  \fill[paleblue] (67,0) rectangle (103,36);
  \foreach \x in {3,12,21,30,39}
    \draw[teal!30,line width=.3pt] (\x,0) -- (\x,36);
  \foreach \y in {0,9,18,27,36}
    \draw[teal!30,line width=.3pt] (3,\y) -- (39,\y);
  \draw[teal,line width=.65pt] (3,0) rectangle (39,36);
  \draw[teal,line width=.65pt] (67,0) rectangle (103,36);
  \foreach \x in {7.5,16.5,25.5,34.5}
    \foreach \y in {4.5,13.5,22.5,31.5}
      \fill[brand] (\x,\y) circle (1.1);
  \draw[brand,line width=1pt] (85,18) -- (103,18) -- (103,36) -- cycle;
  \foreach \x/\y in {85/18,88/18,91/18,94/18,97/18,100/18,103/18,
    103/21,103/24,103/27,103/30,103/33,103/36,
    100/33,97/30,94/27,91/24,88/21}
      \fill[brand] (\x,\y) circle (.85);
  \node[fill=paleblue,inner sep=1pt,text=brand,anchor=north east] at (84,17) {$\Gamma$};
  \node[text=brand,anchor=west] at (105,18) {X};
  \node[text=brand,anchor=south west] at (104,35) {M};
  \node[text=ink] at (21,-7) {覆盖整个区域};
  \node[text=ink] at (85,-7) {沿选定方向采样};
\end{tikzpicture}
\caption{$k$点网格与能带路径的区别（二维示意）。前者用于区域积分或统计取样，后者展示选定路径上的变化。图中以正方形区域和三角形路径说明两种取样方式。}
\label{fig:k-sampling}
\end{figure}


這三類取樣可以在同一專案中並存。例如，一條週期計算流程可以先以 $2\times2\times2$ 網格演示 SCF，再沿推薦高對稱路徑計算能帶，或以 $4\times4\times4$ 網格計算總態密度。後續任務繼承所屬流程的自洽密度，使用各自的取樣方式。本書的矽 DOS 入門流程使用八原子晶胞，能帶流程使用二原子標準原胞，兩者各自完成 SCF。第\ref{ch:results}章還介紹二原子原胞專案中儲存的 DOS 數值與重繪方法；閱讀時應核對資料所屬的晶胞。

能帶圖上的一條平滑曲線並不能說明整個布里淵區積分已經準確。相反，一張很密的 DOS 網格也不會讓一條僅含兩個端點的能帶路徑自動變得平滑。選擇引數時應先問清楚當前節點要完成哪一類任務。

\subsection{高對稱路徑的含義}

對稱性分析根據晶格和原子排列推薦高對稱點與路徑。路徑上的標籤，如 $\Gamma$、X、L 等，指向倒空間中特定的位置；相同標籤在不同晶系中並不總有相同座標。

改變晶胞表示、原子排列或對稱性容差，可能影響推薦路徑。比較兩張能帶圖時，除了泛函和能量零點，還應檢查它們是否採用相同晶胞和同一路徑。不能只按橫軸標籤相似就把曲線逐點相減。

本書使用 \ui{Symmetry \& K-Path} 生成路徑，在 \ui{Band Calculation} 中用 \mbox{\ui{Segment Points}} 設定每段取樣點數。這一數值包含端點。初次執行可以用較少點數觀察整體結構；需要平滑圖形或精細比較曲率時，再增加取樣。

\section{設定電子佔據}

\subsection{電子佔據與費米能}
\index{Feimineng@費米能}

計算得到許多電子態後，還需要決定有多少電子佔據各個態。費米能與電子總數和佔據方式有關，是閱讀能帶和 DOS 時常用的參考。金屬中費米能附近有部分佔據的態，絕緣體和半導體則在佔據態與未佔據態之間可能存在能隙。

數值計算中的費米能不是各個體系都可以直接比較的絕對真空能級。本書能帶圖把費米能移到零點，便於辨認其附近的電子態。如果要比較功函式、能帶對齊或不同體系之間的絕對能級，需要進一步建立一致的能量參照。

\subsection{展寬與電子溫度}
\index{Zhankuan@展寬與電子佔據}

\ui{Smearing} 提供 fermi-dirac、cold、gaussian 和 mp。展寬方法使佔據在能量上較平滑，常用於改善金屬或接近簡併態體系的布里淵區積分和收斂。\ui{Temperature (K)} 給出相應的電子溫度引數。

電子溫度引數不等於給原子加熱，也不意味著計算已經包含某一溫度下的原子振動。一個在 $300$ K 電子佔據下執行的靜態 SCF，原子仍然可以停在固定位置。需要有限溫度結構性質時，應選擇與研究問題相符的方法。

改變展寬和溫度會影響佔據及相應的能量量。比較不同結果時，應說明使用的是總能量、自由能還是其他經過處理的量，並保持定義一致。對於初學教程，先按給定條件執行，再單獨研究展寬變化，閱讀結果會清楚得多。

\subsection{額外能帶}
\index{Ewai nengdai@額外能帶}

電子數決定至少需要多少佔據態，但能帶和 DOS 往往還要檢視費米能以上的空態。\ui{Extra bands} 增加求解的未佔據能帶數量。它不是給體系增加電子，也不會直接改變原子的價態。

若圖中只有佔據態而缺少所關心的導帶區域，應檢查額外能帶數，而不是隻把縱軸範圍拉得更大。座標軸只能顯示已有資料，不能生成沒有計算的電子態。本書矽教程設定額外 $4$ 條能帶，作為觀察低能空態的起點；更高能區需要相應增加。

\section{自旋與磁性}
\index{Zixuan@自旋}

\ui{Nspin} 提供 $1$、$2$、$4$ 等設定，分別對應不同的自旋處理。本書的非磁性示例使用 $1$。存在未配對電子、磁有序或自旋軌道耦合時，應根據實際物理問題選擇適當模式，並檢查相關高階設定。

自旋設定不是一個僅影響速度的開關。氫原子與氫分子具有不同的電子佔據情況；若計算分子解離能，原子參考態的自旋處理會直接影響結果。過渡金屬體系還可能存在多個自洽解，初始磁性和幾何結構都會影響最後得到的狀態。

本書的基礎案例不使用磁性、非共線或自旋軌道耦合。學習時可以先把非磁性工作流掌握，再為一個明確的磁性問題建立新專案。這樣更容易區分結構、取樣和自旋各自造成的變化。

\section{讓自洽計算收斂}

\subsection{對角化方法}

\ui{DiagMethod} 控制電子態的數值求解方法。\ui{davidson\_qe}、\ui{lobpcg} 等方法針對同一電子結構問題，採用不同的迭代策略。充分收斂時，它們應給出一致的物理結果，主要區別在計算效率和求解過程。

比較演算法時，應使用相同輸入和目標精度，以各自收斂後的結果判斷一致性。本書的 SCF 例使用預設方法，後續 NSCF 例保留其預設求解設定。

\subsection{最大步數與收斂精度}

\ui{Max SCF steps} 是允許的最大迭代次數。把它從 $100$ 改為 $200$，只給求解器更多機會，不會自動把精度提高一倍。真正的收斂判斷由誤差和相應閾值決定。到達最大步數但未滿足收斂條件，應視為需要繼續處理的計算。

基礎 SCF 面板提供方法、最大步數和力等選項。密度閾值與混合引數屬於進一步的數值設定，指令碼介面可以使用 \texttt{options.SCF} 配置。本書的基礎教程沿用預設密度閾值 $10^{-8}$。需要調整密度閾值與混合引數時，可在完整MATLAB\regmark{}中使用這一指令碼介面。


設一次電子迭代的輸入與輸出密度之差為
\begin{equation}
  \Delta\rho=\rho_{\mathrm{out}}-\rho_{\mathrm{in}},
  \qquad e_{\mathrm{SCF}}<\varepsilon_{\mathrm{SCF}}.
  \label{eq:scf-stopping}
\end{equation}
左式說明比較的是哪兩份密度；右式說明所用誤差指標必須低於閾值。預設密度混合計算把$\Delta\rho$在倒空間中作加權比較，得到$e_{\mathrm{SCF}}$，並與預設的$10^{-8}$比較。日誌雖使用\file{Rel Rho Err}這一名稱，其值不能直接讀成百分比，也不是兩份密度之差的範數除以輸入密度的範數。判斷本書例項是否收斂，應將這一指標與對應閾值比較，並檢視最終狀態。

\subsection{密度混合的作用}
\index{Midu hunhe@密度混合}

一次求解得到的新密度未必是下一次最好的起點。密度混合把當前結果和歷史資訊組合起來，減小來回振盪並加快收斂。不同體系對混合的響應不同，金屬、大晶胞或磁性體系可能需要更仔細的設定。

若誤差持續在兩個水平之間擺動，可以先檢查結構和佔據，再考慮混合和取樣。增加最大步數可以讓一個已經穩定下降的序列繼續走完，卻未必能解決長期振盪。處理這類問題時，保留完整日誌比只儲存最後一行報錯更有幫助。

\section{幾何最佳化中的兩個精度}
\index{Jihe youhua@幾何最佳化}

幾何最佳化要求每一個候選結構都先完成電子計算，然後根據力更新位置。電子密度收斂精度決定每一步能量和力的可靠程度，幾何最佳化閾值則決定何時認為原子位置已經足夠穩定。


原子$I$的力是總能量對該原子位置的負梯度：
\begin{equation}
  \mathbf F_I=-\nabla_{\mathbf R_I}E,
  \qquad F_{I\alpha}=-\frac{\partial E}{\partial R_{I\alpha}}
  \quad(\alpha=x,y,z).
  \label{eq:atomic-force}
\end{equation}
在計算採用的原子單位中，力以Ha/Bohr表示。力指向能量下降的區域性方向；總能量已經變化很小，並不自動意味著每個原子的力都很小。最佳化記錄還可能列出最大原子力、均方根力或整體梯度範數，比較時應先確認列名與定義。

\ui{RelaxTolerance} 與 \ui{MaxRelaxSteps} 分別控制幾何最佳化的停止尺度和最大工作量。本書氫分子最佳化例明確輸入 $10^{-4}$ 和 $100$。圖形面板和指令碼介面的預設值可以不同，記錄時應寫出實際採用的條件。

最終結果應同時檢查幾何最佳化是否收斂、最後一次電子計算是否收斂，以及最後的原子力是否符合研究所需。只看到鍵長不再明顯變化，仍不足以判斷結構已經達到穩定位置。

\section{選擇雜化交換演算法}

在Build Molecule的\ui{Functional}中選擇HSE06後，計算會包含雜化泛函的交換處理。SCF節點高階選項中的\ui{ExxMethod}用於選擇這一處理的演算法，可見選項有default、normal、ace、qrcp和kmeans。泛函決定要使用的近似，交換演算法決定怎樣組織其中的計算；它們不是可以互相替代的兩個下拉框。

初次使用時先保留default。需要比較演算法時，在一個已經有基態計算記錄的小體系上另存專案，只改變ExxMethod，核對總能量、收斂狀態和用時。比較期間保持泛函、模型尺寸和取樣一致。使用PBE等半局域泛函時，保持default即可。

\section{區分基態與激發態}

SCF的軌道本徵值、兩個軌道的能量差和多電子體系的激發能是不同的量。基態計算已經收斂，只說明該基態問題得到了解答；希望研究激發時，還要選擇響應方法、佔據軌道和空軌道組成的計算空間。

TDDFT任務中的\ui{nvbands}和\ui{ncbands}分別是參與響應計算的佔據與空軌道數，\ui{nroots}是要求的激發態數。它們都不是原子數，也不是k點數。首先在Build Molecule中求出足夠的額外能帶，再在TDDFT裡選擇其中的空軌道。以一個佔據軌道和四個空軌道為例，最多隻有四個這樣的軌道對，不能據此要求十個激發根。

\ui{Use TDA}使用Tamm--Dancoff近似；\ui{Use IPA}和\ui{Use RPA}用於不同響應近似的比較。學習時先使用一套明確的設定，再單獨改變近似觀察差異。TDDFT的輸入準備、適用範圍和引數示例見第\ref{ch:workflows}章\ref{sec:tddft-workflow}節。吸收峰需由相應的響應計算得到。


\section{為一個新體系選擇引數}

開始新專案時，先明確研究目標，例如比較吸附位置的相對能量、觀察應變對能帶的影響，或尋找分子的平衡結構。建立能夠執行的小規模設定，核對結構、元素和電子數，再圍繞目標物理量逐項提高精度。保留一份基準結果，每次改變一個條件，後續比較才有清楚的依據。

隨專案儲存簡短的條件表，記錄結構來源與處理、晶胞、泛函、贗勢、截斷能、k點、佔據、自旋及兩類收斂條件。報告列出影響結論的主要資訊，完整工作流和輸入檔案留在專案資料中，便於以後理解和重複計算。
